\chapter{The History Monad}
\label{ch:histmonad}

\section{Keeping the past}

A rewrite theory that discards its past is irreversible precisely because a configuration
may have many predecessors. Recording which step was taken restores a unique one. That is
the whole idea, and packaging it gives the mirror image of Chapter~\ref{ch:costmonad}.

\begin{definition}[The history endofunctor]
$\Hist$ sends a continued interactive GSLT $G$ to the theory whose configurations are
pairs of a configuration of $G$ and a word in the free monoid of rewrite events, with
each rewrite additionally appending its own event to the word.
\end{definition}

$\Hist$ is the writer monad over the free monoid of rewrite events. Its unit attaches
the empty history; its multiplication concatenates.

\begin{proposition}
\label{ob:prop:hist-adjunction}
$(\Hist, \eta, \mu)$ is a monad on $\catciGSLT$, and the evident functor installing the log
is left adjoint to the functor stripping it, with induced monad $\Hist$.
\end{proposition}

\section{The free history is the reversible tree}

\begin{proposition}[Universal cover]
\label{ob:prop:univcover}
The free history makes the graph of computation paths out of a term into a \emph{tree}:
it is the universal cover of the reduction graph, and reversibility is exactly its
unique-parent property.
\end{proposition}

This is the structural content of the construction and it explains the erasures. An
erasure of history is a monad quotient, and these form a lattice between the reversible
cover and the irreversible base. The injective erasures compress the log without folding
the tree; the identifying ones fold it, closing loops so that paths revisit states. Such
an intermediate object is a quotient of the reduction graph's fundamental group by a
chosen set of loops --- an intermediate cover --- with the order-forgetting
(Mazurkiewicz) erasure folding exactly the commuting concurrent steps.

\begin{remark}[Where this is used]
\label{ob:rem:hist-forward}
Chapter~\ref{sec:reversibility} is this construction under a physical reading.
The reversible envelope $\GSLT^\dagger$ built there is $\Hist\GSLT$ with the
backward rules written down, and the reason writing them down costs nothing is
Proposition~\ref{ob:prop:univcover}: on a tree the inverse of a step is
determined, not chosen.
The lattice of erasures is what that chapter's physics needs and does not have a
name for --- an intermediate cover is a coarse-graining, and which loops one is
permitted to close is which histories one can afford to forget.
\end{remark}

\section{The ledger law}

Underneath the covering-space picture sits an invariant needing no cohomology. Write
$\sigma$ for the fuel remaining and $\kappa$ for the fuel recorded as spent. Then
\[
  \sigma + \kappa = \sigma_0 ,
\]
initial fuel equals fuel remaining plus fuel recorded spent: a potential-plus-kinetic
invariant of a cost-accounted, history-keeping run. Its one-way breakage under erasure is
Landauer's principle~\cite{Landauer} --- the history one may erase for free is exactly
the redundant, state-recoverable part.

\section{The duality}

The two constructions are ledgers of the same time axis, read in opposite directions.

\begin{center}
\begin{tabular}{@{}lll@{}}
\toprule
& \textbf{Cost $\Cost$} & \textbf{History $\Hist$} \\
\midrule
adjoined structure & token stack & event word \\
monoid             & free (list), consumed & free (list), accumulated \\
behaviour along $\longrightarrow$ & drains & accrues \\
records            & what remains & what was spent \\
free object        & fully funded run & universal cover \\
erasure lattice    & grades of funding & subgroups of $\pi_1$ \\
physical reading   & bounded causality & recorded causality \\
\bottomrule
\end{tabular}
\end{center}

A GSLT gives causality. $\Cost$ gives \emph{bounded} causality. $\Hist$ gives the ability
to \emph{record} bounded causality. Running $\Cost$ on an $\Hist$-equipped model gives the
cost of making the recording, and the levels are worth keeping distinct rather than
collapsing: the stacking $\Cost\Hist\Cost\Hist\cdots$ is meaningful, and its fixed point
is an object of interest we do not pursue here.

\begin{remark}[Internalisation, again]
As in Remark~\ref{ob:rem:internalise-cost}, a sufficiently expressive base theory can encode
the logging apparatus into itself, performing the record-keeping with its own
computation. The same caveat applies: this is a different erasure from forgetting, and we
note rather than develop it.
\end{remark}
